\section{Findings: Thematic Analysis}
\label{sec:ch4-findings-thematic}
We present eight themes (T1–T8) developed through a bottom-up reflexive thematic analysis~\cite{braun2023thematic, byrne2021worked}. Each line of interview data was open-coded, with codes iteratively grouped into themes. We include illustrative quotes to ground each theme in expert insight and summarise them in \autoref{tab:themes-overview}. 

\paragraph{T1 - Ecological Authenticity and the Essence of Nature in Building Design:}
\label{theme:t1}
Our analysis reveals the need to re-evaluate biophilic design philosophy~\cite{wilson1986biophilia} in contemporary building practices. Experts critiqued the superficial adoption of biophilic design and the tendency to treat nature as a static background (e.g. through use of artificial plants), reducing it to aesthetic elements devoid of ecological interaction or ecological value. One expert noted:  \textit{"When nature is treated as an object, we miss the opportunity to learn from and collaborate with it, instead imposing our own limited understanding."} Others raised concerns about building rating systems like \textit{WELL certifications}\sidenote{\url{https://standard.wellcertified.com/well}}, arguing that they risk reducing biophilic design to checklists, inadvertently fostering ``green-washing''. 

\begin{mainfigure}[tbp]
  \centering
  \captionsetup[subfigure]{
    justification=centering,
    singlelinecheck=false,
    font=footnotesize
  }

  \begin{subfigure}[t]{.32\linewidth}
    \centering
    \includegraphics[height=4.6cm,keepaspectratio]{images/dress.jpg}
    \caption{Iris van Herpen's \emph{Nautiloid}.}
    \label{fig:t2-fashion}
  \end{subfigure}%
  \hfill
  \begin{subfigure}[t]{.58\linewidth}
    \centering
    \includegraphics[height=4.6cm,keepaspectratio]{images/daan.jpg}
    \caption{Daan Roosegaarde's \emph{Waterlicht}.}
    \label{fig:t2-multisensory}
  \end{subfigure}

  \caption{Examples of biophilic expression through art: (a) Studio Iris van Herpen's creation inspired by aquatic life forms. Image credit: all rights reserved by Iris van Herpen BV, Iris van Herpen. Accessed 14 May 2025. (b) \emph{Waterlicht}, a project showcasing life on a flooded Earth. Image credit: Daan Roosegaarde, \protect\url{https://www.studioroosegaarde.net}. Accessed 14 May 2025.}
  \Description{Two images illustrating artistic interpretations of biophilia.}
  \label{fig:t2-examples}
\end{mainfigure}


In exploring the essence of nature in building design, experts highlighted a non-interventionist mindset for achieving ecological authenticity (discussed further in \autoref{sec:ch4-design-opportunities}). The concept of ``\emph{third landscapes}'' where \textit{space expresses neither power nor submission to power..."}, inspired by \citet{clement2004manifeste}, was mentioned as a model for enabling ecological processes to emerge naturally---designing for conditions rather than dictating outcomes. One expert described this as: \textit{"Instead of dictating (e.g. planting) designs, we should focus on creating foundational conditions like soil and shade to let nature respond and evolve."} These reflections raise questions about the role and representation of nature in design: What constitutes the essence of nature in built environments, and how can we reconcile its authenticity with the functional and experiential expectations of intelligent indoor spaces, while respecting its authenticity?


\paragraph{T2 - Art as an Expression for Biophilic Experiences:}
\label{theme:t2}

Biophilic art is well documented~\cite{miles2010representing, brady2019aesthetics} in HCI through examples such as robotic flowers~\cite{chooi2022symbiosis} and vapour-based spatial partitions~\cite{deng2019diffusive}. Experts highlighted (emerging) technology as an ``artistic material for extending natural behaviours'' and enabling participatory design approaches. Biophilic art must embody ``response'' to reflect nature’s essence (tying to the previous theme of ecological authenticity):  \textit{"...An important layer is reaction, so whenever a human does something, the sensory nature experience reacts. So like grasses in a landscape where you run through with your hand, and then lights start to glow. It’s not natural in the way we know it, but it would give, I think, a very natural response."} Experts suggested that this response could take the form of \emph{soft fascination}---a quality that gently captures attention without overwhelming the mind (explored in HCI~\cite{basu2019attention, sabinson2024every})---creating moments of distraction or relief.


Digitally mediated biophilic art can act as a participatory medium, encouraging occupants to engage with and shape their surroundings, as well as develop emotional and cultural connections to nature:  \textit{"Artistic dimensions can mediate between human-made environments and natural inspirations, crafting spaces that resonate emotionally with users."} This could be through art that spatializes human interaction with nature (e.g. ~\cite{marchart1998art}), such as interactive projections or augmented reality experiences~\cite{shepard2011sentient}. Studio Roosegaarde’s \textit{Waterlicht} installation, which uses light and reflective lenses to simulate virtually flooded environments\sidenote{\url{https://www.studioroosegaarde.net/project/waterlicht}}, was shared as an example. Crafting and 3D printing were highlighted as another participatory medium for tying cultural narratives with nature. Iris van Herpen’s \textit{Architectonics} demonstrates coexistence between humans and aquatic environments through future floating cities fashion\sidenote{\url{https://www.irisvanherpen.com/collections/architectonics}}.



\paragraph{T3 - Biophilic Design for Ecological Empathy:}
\label{theme:t3}
Biophilic design holds the potential to foster ecological awareness and empathy in urbanised societies where disconnection from nature is growing~\cite{li2022re,wang2022nature}. Experts described how interactions with nature influence values like kindness and respect, encouraging empathic behaviours. However, two limitations constrain current approaches.

First, raising awareness without cultivating a lasting commitment to ecological preservation is insufficient. Experts suggested integrating participatory stewardship, where individuals care for natural elements to build a sense of responsibility. Central to this is \emph{``cues for care''} (previously put forth by ~\citet{nassauer1995messy})---design that signals value and the need for maintenance. One expert described a project using QR codes and sensors to trigger pop-ups with information about nearby plants: \textit{"We're adding a technical layer to inform people about what is happening there...At [company name], we wanted to create a sort of walking green route where once they're (employees) near a sensor, they get a pop-up about the plant or the tree so they can learn more about it..."} 

Second, buildings fail to make invisible ecological processes like species health or inter-species interaction perceivable to occupants (e.g. the MicroAquarium installation that reveals live micro-organism activity~\cite{lee2020microaquarium}). These interactions often unfold at temporal scales that are not perceptible to humans, making them harder to translate into meaningful interactions. Experts argued that making such processes visible can reinforce a sense of interconnection: \textit{"By designing buildings that show people how to coexist with other species, we can remind users of their interconnectedness with the living world."} While indoor biotopes are a known strategy~\cite{kerimova2012biotope}, experts felt they often prioritise aesthetics over fostering curiosity and reflection. Such an emphasis limits the ability to inspire empathy and a sustained commitment to preservation.


\paragraph{T4 - Multisensory and Multidimensional Biophilic Design:}
\label{theme:t4}
In literature, we found that visual elements like greenery, views, and patterns are common~\cite{beatley2016handbook, lefosse2023biophilia}. This visual focus has been critiqued, notably by Pallasmaa in \emph{The Eyes of the Skin}. Addressing this, one expert remarked: \textit{"Using visual is an easier way to experience biophilia, and we found that most biophilic designs heavily relied on this sense."} While foundational, visual design alone fails to capture the full richness of how humans interact with nature. Experts stressed that nature inherently presents \textit{layered} sensory experiences, such as visual, auditory, tactile, and olfactory, creating interactions that are often absent indoors. The kinaesthetic sense---our awareness of bodily movement---was emphasised. Experts spoke of designing for kinaesthetic sense in two ways: \textit{movement (with)in nature} (e.g. flowing water, shifting light) and \textit{movement towards nature} (e.g. spatial layouts that invite exploration). One shared:  \textit{"It’s not enough to see trees through a window---design should invite people to go out, move, and interact with the environment."}

Experts described biophilic experiences as multidimensional, offering three key experiential examples: First, nature operates on varying \emph{scales}, from the vastness of landscapes that inspire awe and reflection, to the intimacy of small details that invite personal connections. As one expert explained: \textit{"You are in contact with nature... As humans, you start to reflect on your own ego and not really want to control everything, and you feel a kind of relief."} Second, \emph{temporal rhythms} such as tides, light shifts, or seasonal cycles were also viewed as essential yet rarely integrated indoors. One participant reflected: \textit{"In nature, time is dynamic---you see tides coming in and out, shadows shifting... There’s a rhythm to it that’s deeply tied to our perception of being in nature."} Finally, \emph{infinity} was described as another overlooked aspect: \textit{"Views of the horizon or infinite space---like transparent designs or reflective surfaces---can create an illusion of endlessness that connects users to the vastness of nature."}. 


\paragraph{T5 - Multi-species Flourishing in Buildings:}
\label{theme:t5}
Historically, architecture has excluded or controlled non-human species---Vitruvius’ ``primitive hut''~\cite{laugier1755essay, topccuouglu2022architecture} symbolises architecture’s functional origins, prioritizing human protection and comfort. Modern practices often perpetuate this trend through impermeable designs that hinder ecological integration. Even within buildings that incorporate natural elements, one expert noted: \textit{“Most of the time we get requested a plant with hydro grains... the plant itself is in a hibernation state. It’s surviving but not really actively doing anything.”} In contrast, multi-species flourishing considers buildings as \emph{opportunities} for ecological interaction and care between humans and non-humans.

This perspective informs design approaches suggested by experts. First, embedding \emph{encounters} with non-human species can foster meaningful interactions. As one expert noted: \textit{"Creating spaces for multi-species social encounters can encourage new experiences for us.... and also shared use and mutual respect for the non-human."} Second, \emph{dedicated spaces} for multi-species flourishing are essential. While pollinator gardens~\cite{wenzel2020urbanization} and insect hotels~\cite{james2019urbanisation} are popular outdoors, indoor spaces must support local life, including during building off-hours. Connectivity is vital; transitional spaces like ecological corridors~\cite{zellmer2022urban} can link fragmented habitats, enabling species and ecosystem interaction. 

Technology was suggested in managing these interactions: \textit{"...a smart cat door can open at certain times for a specific cat. You could extend this to ecological passages... sometimes there may be too many mice, so we close it or we try to balance between different ecological communities. Technology could monitor and try to manage this relationship."} Finally, ``Imagination'' was highlighted as a critical design tool (also discussed by \citet{fernandez2024green}): \textit{"Imagination allows us to see beyond functional limitations and engage with the possibilities of truly integrated ecological systems."}.


\paragraph{T6 - Biophilic Design for Social Justice:}
\label{theme:t6}
Two experts raised the potential of biophilic design as a means of addressing social inequities by fostering inclusive, community-focused environments. Research supports this premise---nature and affluence are linked in urban settings, where areas with abundant greenery tend to be wealthier, while under-resourced areas often lack access to nature~\cite{wolch2014urban}. This disparity extends to indoor spaces as well: \textit{“Biophilic design is frequently seen as a luxury, rarely accessible to those who might benefit most from it.”} We saw that most biophilic projects mentioned by experts, whether from their own work or known examples, were heavily concentrated in corporate developments.

Despite this, biophilic design can support marginalised communities by combining ecological stewardship with social inclusivity~\cite{panlasigui2021biophilia}. First, it can foster a sense of ownership through participatory design, where \emph{collaborative biophilic projects} bridge gaps between professional intent and the lived experiences of diverse users. Such projects can promote economic uplift by integrating urban agriculture or green infrastructure, providing community resources and ecological benefits. Second, biophilic design can help meet the emotional and social needs of transient or vulnerable populations by creating spaces that provide refuge and healing (also shown in existing works~\cite{tidballurgent, chow2020site}). Experts emphasised programming to keep such spaces active and inclusive. As one noted, \textit{“Programming ensures that these spaces don’t just exist as beautiful designs but are used. It’s like creating a modern version of a church---a space for community and solace.”} 

In literature, in ancient temple courtyards, biophilia served collective and inclusive functions~\cite{gurung2015expanding}. The industrial era, however, shifted design priorities toward efficiency and density, often at the cost of natural elements. While biophilic design is re-emerging as a response to ecological and social challenges, its implementation often reflects existing inequalities. In corporate settings, it manifests as immersive features like indoor ponds; in underserved contexts, it is reduced to token greenery. Reclaiming biophilic design’s communal, inclusive legacy may help address these disparities and ensure its benefits reach those who need them most.


\paragraph{T7 - The Body in Biophilic Design:}
\label{theme:t7}

The human body is central to biophilic design---as both observer and participant, it senses and responds to environmental cues: \textit{"...the human body—as a natural sensor... it’s the first means we can use to recreate our inner connection with nature."} This view builds on \citet{shusterman2011somaesthetics}, where bodily scale and orientation guide spatial design. \emph{Somaesthetic design} in HCI extends this perspective, highlighting bodily awareness and the reciprocal relationship between body and environment\cite{hook2008knowing, hook2018designing}.

\emph{Body sensitivity} emerged as a critical design principle, advocating for delicate, nuanced interactions with natural elements, like moss textures or rustling leaves, to foster deep somatic connections. Sensitivity was framed through a feminist lens, addressing gendered differences in somatic needs. Experts discussed how hormonal cycles and bodily rhythms influence indoor experiences (also researched in gender studies~\cite{haselsteiner2021gender, kim2013gender}), and how buildings should respond to these differences. Technology was offered as a means to respond to feminist bodily needs: \textit{"And in an ideal building, it is more related to how our body works and how we feel and how we can be the most happy and productive in that sense. And then technology could analyse your body---body temperature, for example. So for a woman’s body, where are you in your cycle? What do you need?"}. Another highlighted the importance of augmenting this sensitivity to make biophilic experiences accessible: \textit{"You need to be sensitive to first catch the different benefits a biophilic experience can offer. And we have seen that the female body is more sensitive..."}. We discuss this in greater detail in \autoref{sec:ch4-design-opportunities}.

Bodily awareness, biophilia, and architecture were also discussed through what one expert termed the ``living effect'' (briefly explored in HCI~\cite{steer2015growth}), highlighting elements that evoke growth and vitality in artificial spaces. This idea extends to the concept of \emph{bodily flow}, where design elements such as meandering paths or kinetic installations not only encourage physical movement but also create an awareness of the  interaction between the body and its surroundings~\cite{shusterman2011somaesthetics, cheng2022sonification}. By actively guiding movement, such designs enable occupants to experience spaces in a way that feels organic and intuitive, enhancing their connection to both the environment and their own embodied presence. The somatic-biophilic lens thus offers an alternative lens on how biophilic design can support deeper connection to nature and self.

\paragraph{T8 - Biophilic Design for Occupant (Dis)Comfort:}
\label{theme:t8}

Experts highlighted that biophilic design should challenge traditional comfort norms: \textit{"Comfort should not always about ease; sometimes it's about engagement. And I think that nature's unpredictability makes us feel alive, even when it's inconvenient."} In smart buildings, for example, sunlight enriches sensory experiences with light and shadow (called Komorebi~\cite{zhang2020komorebi}), but is often ``corrected'' through automated shading systems~\cite{casini2018active}. Similarly, natural elements perceived as ``messy'' or unpredictable are frequently avoided: \textit{"People think it's (referring to a landscape project) very messy, and they are reluctant to interact with the space. But sanitised versions limit and reduce our connection to the natural world."}

In nature, discomfort is often experienced as \emph{duality}. Walking barefoot on sand can be grounding, yet painful (e.g. from sharp shells). One expert reflected: \textit{"...I am really enjoying being in nature but it could also a lot of time be a bit scary, bring fear...like innate fear when you are confronting with nature."}  Yet, these moments of discomfort can catalyse personal reflection: \textit{"Friction fosters reflection. When we're uncomfortable, we're more likely to pay attention and find meaning in our surroundings. A lot of new reflection actually comes from this friction in contact with nature."} This tension between comfort and discomfort seems inevitable. Designing for biophilic (dis)comfort could involve strategic oscillations between ease and challenge, mimicking natural rhythms: \textit{"The stepping stones (installed in the midst of a footpath) reduce paving...you really have to watch your step, but you feel like you’re passing through a natural overgrown area."} Discomfort could also be framed as a trade-off, where users accept minor inconveniences---such as insects from pollinator plants indoors---for greater ecological benefits. This is discussed in \autoref{sec:ch4-design-opportunities}. 

In productivity-focused spaces, biophilic (dis)comfort raises critical questions of balance and appropriateness. While greenery and lighting enhance focus (e.g.~\cite{aries2005human}), one expert noted: \textit{"There's an appetite for more greenery, especially after the pandemic, to try to make sure that the spaces indoors actually do have the qualities that are a bit more, let's say, friendly and encourage people to come back and work."} How can (dis)comfort align with expectations in workplaces or learning environments? Finally, limiting design to ableist bodies is problematic. Overly wild, unpredictable elements may be inaccessible to individuals with mobility or sensory sensitivities: \textit{"you need to see if everything is accessible. You could have parts like that (referring to a project that designed for exploratory pathways paved with stones) but it needs to be wheelchair accessible. Otherwise, it can limit (someone's) movements."}.

\begingroup
\fontsize{7.3}{8.7}\selectfont
\setlength{\tabcolsep}{2.4pt}
\renewcommand{\arraystretch}{1.04}
\setlength{\LTleft}{0pt}
\setlength{\LTright}{0pt}

\begin{longtable}{@{}%
  >{\RaggedRight\arraybackslash}p{.24\linewidth}
  >{\RaggedRight\arraybackslash}p{.43\linewidth}
  >{\RaggedRight\arraybackslash}p{.23\linewidth}
  >{\RaggedLeft\arraybackslash}p{.05\linewidth}@{}}

\caption{An overview of the eight themes created from a reflexive analysis of the expert interviews. For each theme, we provide the theme name, a brief description, the top three codes associated with the theme, and the number of supporting quotes.}
\label{tab:themes-overview}\\

\toprule
\textbf{Theme} & \textbf{Description} & \textbf{Example Codes (n=3)} & \textbf{Count} \\
\midrule
\endfirsthead

\caption[]{An overview of the eight themes created from a reflexive analysis of the expert interviews, continued.}\\
\toprule
\textbf{Theme} & \textbf{Description} & \textbf{Example Codes (n=3)} & \textbf{Count} \\
\midrule
\endhead

\midrule
\multicolumn{4}{r}{\footnotesize Continued on next page}\\
\endfoot

\bottomrule
\endlastfoot

T1 - Ecological Authenticity and the Essence of Nature in Building Design &
This theme revisits the philosophy of biophilic design as a deep, authentic connection with nature that aligns with its processes, patterns, and non-human agency. However, this philosophy is often diluted or lost in translation during the design process, leading to tokenistic or anthropocentric implementations. The focus is on preserving the integrity of ecosystems by designing in harmony with their natural cycles rather than imposing human-centred control. &
Shallow adoption and Tokenism; Re-framing Nature as Active Participant; Simplified perceptions of nature &
75 \\

\midrule

T2 - Art as an Expression for Biophilic Experiences &
This theme examines our inclination to represent and connect with nature through artistic expression. It discusses how art and technology can combine to enable biophilic experiences to create participatory, restorative, and culturally resonant designs. &
Technology as participatory ``material''; Soft fascination; Natural behaviours through biophilic art &
29 \\

\midrule

T3 - Biophilic Design for Ecological Empathy &
Biophilic design emerges as a medium for fostering environmental awareness, ecological empathy, and sustainable behaviours. By integrating biophilic principles into design, spaces can support active learning, encouraging a sense of stewardship and collective responsibility for the natural world. This theme emphasises how educational biophilia---using interactive, sensory, and participatory approaches---can cultivate biophilic behaviours in occupants. &
Nature awareness as a catalyst for nature respect; Participatory stewardship of nature; Technology for amplifying living elements &
52 \\

\midrule

T4 - Multi-sensory and Multi-dimensional Biophilic Design &
This theme emphasises the need for expanding biophilic design in indoor spaces beyond visual and greenery-centric approaches to include sound, smell, touch, and movement as well as including more complex and layered environmental interactions. &
multi-sensory experience, Experience of sublime, Movement and fluidity in design. &
88 \\

\midrule

T5 - Multi-species Flourishing in Buildings &
This theme envisions biophilic design for supporting environments where diverse species can coexist, thrive, and contribute to ecological resilience. It emphasises the agency of non-human species as active building ``occupants''. The focus is on creating opportunities for non-human species to express their natural behaviours, flourish in human-altered spaces, and contribute to shared ecosystems. &
Living ecosystem design approach, Integrating with inter-species practices, Non-interventionist design approach &
65 \\

\midrule

T6 - Biophilic Design for Social Justice &
This theme examines biophilic design for social inequities by fostering inclusive, participatory, and community-focused environments. It critiques the exclusivity of biophilic design as a ``corporate culture'' and advocates for equitable access to natural, restorative spaces. The focus is on empowering marginalised and vulnerable communities through designs that provide safe spaces for emotional escape and support active social programming. &
Biophilic design as an introverted design privilege, Spaces of injustice, Biophilia for community building &
63 \\

\midrule

T7 - The Body in Biophilic Design &
This theme explores the role of the human body as a site of interaction and perception within biophilic design. Drawing on somaesthetic principles in built environments, which emphasise bodily awareness, sensory engagement, and movement~\cite{shusterman2011somaesthetics}, this theme explores how somaesthetic design can provide an alternative lens to can foster meaningful biophilic experiences, fostering awareness of nature, and of self. &
Human body as a sensor for biophilic experiences, sensitivity as a BD principle, Movement and fluidity for bodily design &
30 \\

\midrule

T8 - Biophilic Design for Occupant (Dis)Comfort &
This theme challenges traditional notions of comfort by creating environments that evoke both comfort and discomfort. It examines how embracing elements of discomfort including unpredictability and even fear can be a meaningful aspect of the nature experience within indoor spaces. By embracing these ``unexpected'' aspects of nature, this theme examines the opportunities and challenges around designing for biophilic (dis)comfort as a means for connecting with nature. &
Friction in contact with nature; Rejection of Uncontrolled Nature; wildness versus usability &
42 \\

\end{longtable}

\Description{A summary table of eight themes derived from reflexive analysis of expert interviews on biophilic design. For each theme, the table provides the theme name, a descriptive summary, three example codes, and the number of supporting quotes. Themes include: (1) Ecological Authenticity, critiquing tokenistic approaches; (2) Art as Biophilic Expression; (3) Ecological Empathy and Stewardship; (4) Multi-sensory and Layered Design; (5) Multi-species Flourishing; (6) Biophilic Design for Social Justice; (7) Somaesthetic and Embodied Biophilia; and (8) Biophilic (Dis)Comfort, which considers the role of unpredictability and discomfort. The quote counts range from 29 to 88, indicating varying emphasis across themes.}

\endgroup